\documentclass[10pt,conference]{IEEEtran}
\usepackage[T1]{fontenc}
\usepackage{times}
\usepackage{amsmath,amssymb}
\usepackage{booktabs}
\usepackage{listings}
\usepackage{xcolor}
\usepackage[hidelinks]{hyperref}
\usepackage{url}
\hypersetup{pdftitle={A Genealogy of Optimizers},pdfauthor={Siddharth Choudhary and Claude (Anthropic)}}

\lstset{
  basicstyle=\ttfamily\footnotesize,
  breaklines=true,
  columns=fullflexible,
  keepspaces=true,
  xleftmargin=1em,
  frame=leftline,
  framesep=6pt,
  rulecolor=\color{gray!50},
  aboveskip=6pt, belowskip=6pt,
}

\newcommand{\sign}{\operatorname{sign}}

\title{A Genealogy of Optimizers}
\author{\IEEEauthorblockN{Siddharth Choudhary and Claude (Anthropic)}}

\begin{document}
\maketitle

\begin{abstract}
This is a guided tour through the optimizers used to train neural networks, from
stochastic gradient descent to the recent matrix-aware methods Muon, Shampoo,
and MuonH. The organizing idea is that each optimizer fixes a specific, namable
failure of the one before it: momentum gives the optimizer a memory, Adam adds
per-parameter learning rates, AdamW corrects how weight decay interacts with
them, the sign-based methods strip the machinery back to its load-bearing core,
and the matrix-aware methods exploit structure that per-parameter methods throw
away. We trace that lineage, and along the way isolate a single recurring axis
--- how each method turns a gradient into a step --- that explains most of the
progress. An interactive version with live simulations is available online.\footnote{\url{https://itzsid.github.io/publications/optimizer-ladder.html}}
\end{abstract}

\section{SGD and the stochastic gradient}
A neural network is a giant pile of numbers, the parameters. Training means
finding the values that make the model's predictions good, measured by a single
scalar called the \emph{loss}: lower is better. So training is a search problem.
Picture the loss as the height of a landscape and every setting of the
parameters as a location on it; training is finding the lowest valley while
wandering blindfolded.

You cannot see the landscape, but at any point you can feel which way the ground
tilts. That tilt is the gradient, and it points uphill, so you step the opposite
way:
\begin{equation}
w \leftarrow w - \eta\, g,
\end{equation}
where $g=\nabla_w L$ and the step size $\eta$ (the learning rate) is the one knob
you really tune --- too small and you crawl, too big and you overshoot.

The \emph{stochastic} part is the trick that makes this practical. Computing the
true gradient means evaluating the loss on every training example before taking a
single step, which is wildly impractical. Instead we compute a noisy estimate of
the gradient on a small random batch: wrong in the details but right on average,
and nearly free. So instead of one careful step on all the data we take thousands
of slightly drunk steps on slivers of it. The noise that looks like a bug is part
of why SGD works at scale --- random kicks help models escape bad regions.

Three problems are already visible. First, even at the optimum the noise never
goes away, so $w$ keeps wobbling. Second, every step is memoryless: if the last
twenty gradients all said ``go right,'' the twenty-first ignores that history.
Third, in real models different parameters have wildly different gradient scales,
and one global learning rate is asking a lot. Momentum fixes the second and
partly the first; RMSProp and Adam tackle the third.

\section{Momentum}
Stretch the loss bowl into a long narrow valley --- steep walls one way, a gentle
slope along the length. This is what real loss landscapes look like: a few
directions are extremely curved and most are nearly flat. Plain SGD lurches
across the steep walls and barely creeps along the valley floor, because its
steps are memoryless.

The fix is to treat the gradient not as the direction to move but as a force on a
velocity that the optimizer carries from step to step:
\begin{equation}
v \leftarrow \beta v + g, \qquad w \leftarrow w - \eta\, v.
\end{equation}
With $\beta$ near $1$ (typically $0.9$), useless directions whose gradients flip
sign cancel in $v$, while the consistent valley direction accumulates --- at
steady state with $\beta=0.9$ the velocity is roughly ten gradients' worth, so
progress along the valley is about ten times faster. Under noise, the independent
noise terms also average toward zero in $v$ while the signal survives.

Momentum fixes the memoryless complaint and partly the noise one, but it does
\emph{not} solve heterogeneous gradient scales: a parameter with a $100\times$
larger gradient still gets a $100\times$ larger step. That is the next target.

\section{Adam: from AdaGrad to RMSProp to Adam}
In a real network the single global learning rate is catastrophic. The bias of a
softmax output and a weight two layers back can have gradient magnitudes
differing by $1000\times$. The field converged on per-parameter learning rates
computed automatically from gradient statistics, in three steps.

\textbf{AdaGrad} keeps a running sum of squared gradients per parameter and
divides the step by its square root, $w \leftarrow w - \eta\, g/\sqrt{G+\epsilon}$
with $G \leftarrow G + g^2$. Loud parameters get shrunk, quiet ones keep a
near-full rate. The catch: $G$ only grows, so after a few thousand steps every
effective learning rate approaches zero and the optimizer stalls.

\textbf{RMSProp} replaces the running sum with an exponential moving average,
$v \leftarrow \beta_2 v + (1-\beta_2) g^2$, so $v$ tracks \emph{recent} magnitude
and plateaus instead of growing forever \cite{rmsprop}. \textbf{Adam} then stacks
the two ideas --- a momentum EMA of the gradient and an RMSProp EMA of the
squared gradient, both bias-corrected \cite{adam}:
\begin{equation}
\begin{aligned}
m &\leftarrow \beta_1 m + (1-\beta_1)\, g, &
v &\leftarrow \beta_2 v + (1-\beta_2)\, g^2, \\
\hat m &= m/(1-\beta_1^t), &
\hat v &= v/(1-\beta_2^t), \\
\end{aligned}
\end{equation}
\begin{equation}
w \leftarrow w - \eta\, \hat m / (\sqrt{\hat v}+\epsilon).
\end{equation}
The defaults $\beta_1=0.9$, $\beta_2=0.999$, $\epsilon=10^{-8}$ have proven nearly
universal. The punchline took the field years to internalize: if gradients are
roughly stationary then $\hat v \approx g^2$ and $\hat m \approx g$, so
$\hat m/\sqrt{\hat v} \approx \sign(g)$. Adam's update is a momentum-smoothed sign
of the gradient: every parameter moves by about $\eta$ regardless of its gradient
magnitude. That is why default Adam works on heterogeneous networks, and why it
became the reflexive choice for nearly a decade.

\section{Adam as a signal-to-noise ratio}
A cleaner reading of $\hat m/\sqrt{\hat v}$ is a per-parameter signal-to-noise
ratio. Write each gradient as $\text{signal}+\text{noise}$ with noise standard
deviation $\sigma$. Then $m \to \text{signal}$ (noise cancels in the mean),
$v \to \text{signal}^2+\sigma^2$ (the variance survives in the second moment), so
\begin{equation}
\frac{m}{\sqrt{v}} \to \frac{\text{signal}}{\sqrt{\text{signal}^2+\sigma^2}},
\end{equation}
exactly the signal-to-noise ratio, bounded in $[-1,1]$ in steady state. Table~\ref{tab:snr}
walks through the regimes: a strong clean gradient gives a full step, a noisy one
a small step, pure noise or oscillation gives nothing, and magnitude alone (rows
A versus F) does not change the answer.

\begin{table}[t]
\centering
\caption{Adam's update as a signal-to-noise ratio. Only the ratio of signal to
magnitude matters, not the magnitude.}
\label{tab:snr}
\small
\begin{tabular}{@{}llrrrr@{}}
\toprule
scenario & signal & $\sigma$ & $m$ & $\sqrt{v}$ & update \\
\midrule
A clean strong   & 0.3 & 0   & 0.3 & 0.3  & 1.00 \\
B moderate noise & 0.3 & 1.0 & 0.3 & 1.04 & 0.29 \\
C weak signal    & 0.1 & 1.0 & 0.1 & 1.00 & 0.10 \\
D pure noise     & 0   & 1.0 & 0   & 1.0  & 0.00 \\
E oscillating    & $\pm1$ & 0 & 0 & 1.0 & 0.00 \\
F loud parameter & 100 & 0   & 100 & 100  & 1.00 \\
\bottomrule
\end{tabular}
\end{table}

One caveat: the bound holds in steady state. During transients --- early
training, or just after a learning-rate change --- the mismatch between the short
$\beta_1$ window and the long $\beta_2$ window can push the update above $1$. This
is part of why gradient clipping is standard for transformer training.

\section{AdamW: the one-line fix}
Adam was published in 2014; AdamW, the version everyone actually uses, came in
2017, and the change is one line \cite{adamw}. Weight decay gently pulls
parameters toward zero each step to improve generalization. For SGD there is an
identity: shrinking the parameters is equivalent to adding $\lambda w$ to the
gradient, so for years people implemented decay by folding $\lambda w$ into the
gradient. That is correct for any optimizer that just multiplies the gradient by
$\eta$ and subtracts.

Adam does not. It divides by $\sqrt{v}$ first, so a folded-in decay term becomes
$\eta\,\lambda w/\sqrt{v}$: parameters with large $v$ (loud gradients) get decayed
\emph{less}, quiet ones \emph{more}. This is exactly backwards --- the parameters
most at risk of overfitting receive the least regularization. The fix is to
decouple decay from the gradient and apply it directly to the weights:
\begin{equation}
w \leftarrow w - \eta\, \hat m/(\sqrt{\hat v}+\epsilon), \qquad
w \leftarrow w - \eta\,\lambda\, w.
\end{equation}
Now every parameter gets the same proportional decay regardless of $v$, and
$\lambda$ is an independent knob. AdamW trained essentially every transformer
between 2017 and roughly 2023 --- the silent default behind GPT-2, GPT-3, BERT,
T5, ViT, and CLIP.

\section{Interlude: SignSGD}
So far the story has been ``track more statistics, use them more cleverly.''
SignSGD asks the rude question: what if we track \emph{less}? Since Adam's update
is approximately $\sign(g)$, why not compute the sign directly \cite{signsgd}:
\begin{equation}
w \leftarrow w - \eta\, \sign(g).
\end{equation}
No $m$, no $v$, no per-parameter scale --- zero optimizer state. The wins are
concrete: AdamW costs two extra values per parameter (roughly $560$\,GB of $m$ and
$v$ for a $70$B model in fp32), while SignSGD costs nothing, and in distributed
training each gradient compresses to one bit, a $32\times$ bandwidth reduction.

But it has two real problems. It carries no magnitude information, so a parameter
near its optimum overshoots by a fixed amount forever. And it is fragile to noise:
the sign of a near-zero gradient buried in noise is essentially random, and
SignSGD takes a full step in that random direction. SignSGD is rarely used as-is,
but it cleanly shows that the \emph{directionally normalized step} --- everything
moves by about $\eta$ --- is doing most of the work that makes Adam great. Add
back a single momentum buffer and most of the brittleness disappears.

\section{Lion: one buffer instead of two}
Lion emerged from a large program search over optimizer update rules and is what
won \cite{lion}. It is substantially simpler than AdamW, uses half the optimizer
state, and is competitive across many large-scale benchmarks (its wins are
task- and tuning-dependent). The rule smooths the gradient and then takes the
sign, with two different mixing rates:
\begin{equation}
\begin{aligned}
\text{update} &= \sign(\beta_1 m + (1-\beta_1) g), \\
w &\leftarrow w - \eta\,(\text{update} + \lambda w), \\
m &\leftarrow \beta_2 m + (1-\beta_2) g,
\end{aligned}
\end{equation}
with $\beta_1=0.9$, $\beta_2=0.99$. It steps using a fast EMA (window $\approx 10$)
but stores a slow one (window $\approx 100$); the asymmetry materially
outperforms using one rate for both. Lion's update is in $\{-1,0,+1\}$ per
parameter --- strict sign, no fractional confidence --- and its per-parameter
normalization comes from the sign function rather than $\sqrt{v}$. What is
intellectually interesting is not that Lion beats AdamW (sometimes it does,
sometimes not) but that it is competitive while throwing away half the machinery.

\section{The hidden invariant: per-parameter normalization}
Looking back from SGD to Lion, a single axis explains most of why each method
outperforms the last: how does a parameter's gradient magnitude affect how far it
moves? Table~\ref{tab:norm} makes it explicit.

\begin{table}[t]
\centering
\caption{Effective step per parameter $i$, and whether it scales with the
gradient magnitude. The adaptive and sign-based methods flatten it.}
\label{tab:norm}
\small
\begin{tabular}{@{}lll@{}}
\toprule
optimizer & effective step & scales with $|g_i|$? \\
\midrule
SGD            & $\eta\, g_i$ & yes, linearly \\
SGD+momentum   & $\eta\, v_i$ & yes, linearly \\
AdamW          & $\eta\, m_i/\sqrt{v_i}\approx\eta\,\sign(g_i)$ & no, flat \\
Lion           & $\eta\,\sign(\beta m_i + (1{-}\beta) g_i)$ & no, flat \\
SignSGD        & $\eta\,\sign(g_i)$ & no, flat \\
\bottomrule
\end{tabular}
\end{table}

This corrects a common error: momentum does \emph{not} solve heterogeneity. It
accumulates gradients over time within each parameter (temporal smoothing) but
preserves per-parameter magnitude. The thing that solved heterogeneity was the
$\sqrt{v}$ divisor or the sign function, both of which produce roughly equal step
sizes per parameter --- spatial normalization across parameters. This reframes
Lion's success: it is not that $\sign(\cdot)$ is smarter than $m/\sqrt{v}$, but
that both reach the same destination, flat per-parameter steps, and Lion gets
there with one buffer and one operation. It also points at the next question: is
per-parameter even the right granularity? Real parameters are organized into
matrices, which have richer structure.

\section{Muon: per-matrix orthogonalization}
Every method so far smuggles in the assumption that the parameters are a flat
list. Adam scales each scalar entry of a weight matrix independently; the fact
that those scalars form a matrix --- with rows, columns, and structure that
survives matrix multiplication --- is thrown away. Muon's bet is that this
structure is exactly the information you need \cite{muon,muon-scalable}.

Any matrix has a singular value decomposition $M=U\Sigma V^\top$, where $\Sigma$
holds the singular values --- how much the matrix stretches space along each
principal direction. Gradient matrices in training have very unequal singular
values: a few directions dominate, most are nearly zero, so a vanilla step makes
progress in one or two directions and barely moves the rest. Muon computes the
momentum-smoothed gradient and then flattens all its singular values to $1$
before stepping:
\begin{equation}
u \leftarrow \beta u + g, \quad \hat u = \operatorname{orth}(u), \quad
W \leftarrow W - \eta\, \hat u,
\end{equation}
where $\operatorname{orth}(\cdot)$ replaces $\Sigma$ with the identity (keeping
$U,V$). Computing this by SVD every step is too slow, so Muon uses a degree-5
Newton--Schulz iteration --- a polynomial in the matrix, applied about five times,
that pushes the singular values toward $1$ using only matrix multiplications. The
standard coefficients are tuned for speed, not accuracy: they leave the singular
values in a band near $1$ rather than exactly at $1$.

Muon applies to hidden weight matrices only; embeddings, biases, and norm
parameters use AdamW alongside it. Practically, Muon set the speed records for
nanoGPT training and, after work showing it scales, was used at the
trillion-parameter scale for Kimi K2 \cite{kimik2}. It also transfers its learning
rate cleanly across model sizes when paired with $\mu$P-style width scaling
\cite{mup}.

\section{The right kind of uniform: directions, not entries}
The normalization story has a subtlety worth isolating. ``Uniform per entry,''
as Adam and SignSGD give, is uniform in the basis of matrix \emph{coordinates}
--- this row, that column --- but that basis is arbitrary. What a weight matrix
\emph{does} is set by its singular \emph{directions}; rotate the coordinate frame
and every entry changes while the function does not move. So the real question is
not whether every entry moved equally, but whether the update pushed equally hard
in every direction it acts on. A per-entry-uniform step can still be wildly
lopsided in direction space, pouring almost all of its effect into one or two
singular directions; when the gradient concentrates, $\sign(\cdot)$ of a
near-rank-one matrix is still near-rank-one, funneling the entire step into a
single direction.

Muon's orthogonalization keeps the gradient's directions but flattens its
singular values, so the update is isotropic: it advances every direction
equally. This is a genuinely high-dimensional effect (in two dimensions, sign
matrices are usually orthogonal and the distinction vanishes). The deeper reason
per-matrix beats per-parameter is basis-independence: per-entry normalization asks
``is each number stepping equally,'' a question whose answer depends on how the
matrix was laid out, while Muon asks ``is each direction stepping equally,'' a
question about the layer's real geometry.

\section{What first-order optimizers can't see}
Every method so far --- including Muon --- uses one piece of information per
parameter: the gradient, which gives the slope but not the \emph{curvature}. Two
parameters with the same gradient but different curvature want very different
steps: the gently curving one can take a big step safely, the sharply curving one
is about to overshoot. For a single parameter, curvature is the second derivative
$H$, and the quadratic Taylor model $L(w+\Delta)\approx L(w)+g\Delta+\tfrac12 H
\Delta^2$ is minimized by Newton's step $\Delta=-g/H$ --- walk until the slope
hits zero.

For $N$ parameters the curvature is the Hessian $H_{ij}=\partial^2 L/\partial w_i
\partial w_j$, an $N\times N$ matrix. Materializing it needs $N$ backward passes;
for modern models with $N\approx 10^9$--$10^{12}$ that is not merely expensive but
physically impossible, so full Newton is dead on arrival. The escape hatch: we
rarely want $H$ itself, only $H v$ for some vector, and that is one extra backward
pass (a Hessian-vector product). The practical second-order story is therefore:
do not materialize $H$, do not invert it, get cheap factored approximations that
match the matrix structure of the network. Shampoo is the field's main
matrix-aware answer.

\section{Shampoo: curvature-aware preconditioning}
Shampoo predates Muon (2018) and was the first matrix-aware optimizer to gain
attention; Muon got most of the practical traction \cite{shampoo}. Shampoo
approximates the full Hessian with two matrices per layer, one for how rows
interact and one for columns:
\begin{equation}
\begin{aligned}
L &\leftarrow \beta L + (1-\beta)\, G G^\top, \\
R &\leftarrow \beta R + (1-\beta)\, G^\top G, \\
W &\leftarrow W - \eta\, L^{-1/4}\, G\, R^{-1/4}.
\end{aligned}
\end{equation}
The $-1/4$ exponents come from a Kronecker-factorization argument: the Hessian is
approximately $L\otimes R$, so its inverse square root applied to $G$ multiplies
by $L^{-1/4}$ on the left and $R^{-1/4}$ on the right. (The original 2018 version
accumulated $L,R$ as un-decayed sums, exactly the AdaGrad form; modern variants
use the EMA shown here \cite{soap}.)

\textbf{Muon versus Shampoo.} The two answer different questions about the same
gradient matrix. Muon asks which directions the gradient spans and steps equally
along all of them, whitening the update's singular values; it is purely
first-order and cheap every step. Shampoo asks how curved the loss is along rows
and columns and rescales by that estimated curvature; it is second-order-flavored
and pays for it with eigendecompositions. At small and moderate scale the two are
about tied, and Shampoo sometimes edges ahead. The gap opens at large scale for
reasons unrelated to the update rule: Muon transfers its learning rate cleanly
under $\mu$P, Shampoo's matrix roots are expensive enough that the preconditioner
is recomputed only every few hundred steps (running on stale curvature), and it
has more knobs. SOAP --- Shampoo's preconditioner with Adam in its eigenbasis ---
folds the two together and often beats both \cite{soap}.

One fairness point that bites in practice: Shampoo is only competitive when it
preconditions a \emph{momentum-smoothed} gradient. Feed its curvature factors a
raw, noisy minibatch gradient and the $L^{-1/4}, R^{-1/4}$ terms amplify the noise
--- a near-zero curvature estimate becomes a large multiplier. Muon's
noise-robustness comes from the same place, its own momentum buffer.

\section{MuonH: pinning the weight norm}
Every method so far concerns the update \emph{direction}. MuonH instead asks how
to control the \emph{magnitude} of the weights \cite{muonh}. Modern transformers
wrap each weight matrix in RMSNorm, which is scale-invariant: scaling the input by
any positive factor leaves the output unchanged, so the weight matrix's magnitude
does not affect the function the layer computes --- only its direction does. Yet
weights still grow during training (roughly as $\sqrt{t}$ without decay), which
inflates the parameter scale relative to the step. AdamW handles this with
decoupled decay; MuonH handles it more radically by \emph{fixing} the magnitude.

Pick a target radius $R=\|W_0\|_F$ and, after every step, project the weight
matrix back to that radius:
\begin{equation}
\begin{aligned}
U_t &= -\eta\, R\, \operatorname{Normalize}(\text{update}), \\
W_{t+1} &= R\, \operatorname{Normalize}(W_t + U_t).
\end{aligned}
\end{equation}
The weight lives on a hypersphere of radius $R$ for all of training, and $\eta$
becomes the relative step size in units of weight norm. Geometrically, the
retraction cancels any radial component of the update, so only tangential motion
survives.

MuonH does not uniformly beat MuonW (Muon with ordinary weight decay). At lower
learning rates MuonH wins; at higher rates MuonW pulls ahead, because letting the
norm grow provides an implicit annealing schedule. There is also a failure mode at
scale: the retraction multiplies the matrix by a single scalar, shrinking all
singular values proportionally, so the trailing singular values that Muon
patiently grows get repeatedly pulled back down and the spectrum collapses into a
few dominant modes --- the ``spectral squeeze.'' Weight magnitude has nonetheless
emerged as a third axis of optimizer design, orthogonal to update direction and
curvature-awareness.

\section{Coda: the through-line}
SGD took small noisy steps along the gradient. Momentum gave the optimizer a
memory. Adam added per-parameter learning rates that we then read as a
signal-to-noise ratio. AdamW noticed the obvious way to add weight decay silently
broke it. SignSGD and Lion showed the per-parameter normalized step was the active
ingredient all along. Muon jumped from per-parameter to per-matrix,
orthogonalizing each update; Shampoo tracked actual curvature per row and column;
MuonH opened a third axis, controlling weight magnitude directly. Each method
either asked what the previous one was failing to use, or what it was doing
redundantly that could be thrown away --- and both directions produced wins. The
matrix-aware methods are new and evolving, but the through-line should hold: every
optimizer fixes a specific failure of the one before it, and noticing the failure
clearly is the hard part.

\begin{thebibliography}{1}
\footnotesize
\bibitem{rmsprop} T. Tieleman and G. Hinton, ``Lecture 6.5 --- RMSProp,''
  Coursera: Neural Networks for Machine Learning, 2012.
\bibitem{adam} D. P. Kingma and J. Ba, ``Adam: A Method for Stochastic
  Optimization,'' ICLR, 2015. arXiv:1412.6980.
\bibitem{adamw} I. Loshchilov and F. Hutter, ``Decoupled Weight Decay
  Regularization,'' ICLR, 2019. arXiv:1711.05101.
\bibitem{signsgd} J. Bernstein et al., ``signSGD: Compressed Optimisation for
  Non-Convex Problems,'' ICML, 2018. arXiv:1802.04434.
\bibitem{lion} X. Chen et al., ``Symbolic Discovery of Optimization Algorithms
  (Lion),'' NeurIPS, 2023. arXiv:2302.06675.
\bibitem{shampoo} V. Gupta, T. Koren, and Y. Singer, ``Shampoo: Preconditioned
  Stochastic Tensor Optimization,'' ICML, 2018. arXiv:1802.09568.
\bibitem{soap} N. Vyas et al., ``SOAP: Improving and Stabilizing Shampoo using
  Adam,'' 2024. arXiv:2409.11321.
\bibitem{muon} K. Jordan, ``Muon: An Optimizer for Hidden Layers in Neural
  Networks,'' 2024. \url{https://kellerjordan.github.io/posts/muon/}.
\bibitem{muon-scalable} J. Liu et al., ``Muon is Scalable for LLM Training,''
  2025. arXiv:2502.16982.
\bibitem{kimik2} Kimi Team, ``Kimi K2: Open Agentic Intelligence,'' 2025.
  arXiv:2507.20534.
\bibitem{mup} G. Yang et al., ``Tensor Programs V: Tuning Large Neural Networks
  via Zero-Shot Hyperparameter Transfer,'' NeurIPS, 2021. arXiv:2203.03466.
\bibitem{muonh} L. Ren et al., ``Rethinking Language Model Scaling under
  Transferable Hypersphere Optimization,'' 2026. arXiv:2603.28743.
\end{thebibliography}

\end{document}
